% This is a file that modifies the file naming conventions in INRSTeX
% for Berkeley Unix ... It works in conjunction with a command file that
% renames old versions of files produced by INRSTeX. The convention is 
% to add "~" to the end of the filename. This is the same indication that
% EMACS uses for old versions.  Note that there is no concern with file
% name length in Berkeley Unix.  The "~" is not sacrosanct. 

% if a job is aborted, the old files may be made current (in csh) with 
%  foreach i (<jobname>.*~)
%    mv $i `basename $i \~`
%  end

% There are several places in INRSTeX that require changes. 

\catcode`\~=11

%  ------ auto.tex -----

% \inputtagfiles={\inputwithcheck {\jobname.tag}}
\inputtagfiles={\inputwithcheck {\jobname.tag~}} % will have been renamed


% -------- tocform.tex ------

%\def\c@k#1{\csname if#1list\endcsname \jobname.#1;-1 \else \jobname.#1 \fi}
%\def\newlistfilename#1{\ifnotdefined \jobname.#1 \else \c@k{#1} \fi}
             % VAX/VMS \c@k is necessary because of the way TeX expands 

\def\c@k#1{\csname if#1list\endcsname \jobname.#1~ \else \jobname.#1 \fi}
\def\newlistfilename#1{\ifnotdefined \jobname.#1 \else \c@k{#1} \fi}
                      % \c@k is necessary because of the way TeX expands 

% ----------- cite.tex ---------

% This does not need changing as long as the citation list ALWAYS follows
% the citations. If this is not the case, then an old version will have to
% used. This might lead to other problems later. 

% \citetagfilename = {\jobname.ctg} %compatible Ugh!

\catcode `\~=\active
